% SPDX-License-Identifier: Apache-2.0
%
% Razboj, a minimal GPU in TxHDL. Built by //docs:razboj. Every listing
% is included from the tree, the picture is what Razboj drew when the
% build ran it, and the waveform is that run's.
\documentclass[journal,9pt]{IEEEtran}

\newcommand{\listingsize}{\fontsize{6}{7}\selectfont}
% SPDX-License-Identifier: Apache-2.0
%
% The house preamble, shared by every document in this directory. Loaded
% after \documentclass, before \title. Keeping it in one file is what
% stops two articles drifting apart in font, listing style or figure
% style.
% Computer Modern. IEEEtran selects Times (ptm) by default, so the face has
% to be chosen explicitly.
%
% Latin Modern is the Computer Modern variety used here, and the choice is
% forced rather than aesthetic. Under T1 encoding, plain `cmr` has no Type 1
% outlines unless cm-super is installed, and it is not in the pinned
% distribution, so pdflatex falls back to EC bitmaps and every glyph in the
% PDF becomes a Type 3 font. That was measured: the first build produced a
% document with 10 Type 3 fonts and no Type 1. Latin Modern is Computer
% Modern redrawn with a full T1 repertoire and real .pfb outlines, so the
% same design comes out scalable.
\usepackage[T1]{fontenc}
\usepackage[utf8]{inputenc}
\usepackage{lmodern}
% microtype is not in the pinned TeX distribution. emergencystretch alone
% keeps the two-column measure from producing overfull lines.
\setlength{\emergencystretch}{3em}

\usepackage{amsmath}
\usepackage{amssymb}
\usepackage{amsthm}
\usepackage{booktabs}
\usepackage{array}
% enumitem is not in the pinned TeX distribution; the list spacing below
% does what \setlist would have done.
\usepackage{float}
\usepackage{xcolor}
\usepackage{listings}
\usepackage{tikz}
\usetikzlibrary{arrows.meta,positioning,fit,backgrounds,calc}
\usepackage[hidelinks,bookmarks=true]{hyperref}

% Numbered, definition-style box for a rule the reader is meant to apply.
\theoremstyle{definition}
\newtheorem{rulebox}{Rule}
\newtheorem{example}{Example}

% Unnumbered asides.
\theoremstyle{remark}
\newtheorem*{tip}{Tip}
\newtheorem*{pitfall}{Pitfall}

% Tighter lists, in place of enumitem's \setlist.
\setlength{\topsep}{2pt}
\setlength{\itemsep}{1pt}
\setlength{\parsep}{0pt}

\newcommand{\code}[1]{\texttt{#1}}
\newcommand{\term}[1]{\emph{#1}}

% Syntax highlighting. Four roles, distinguishable in colour and still
% distinguishable in greyscale, because an IEEE article gets printed: the
% keyword is the only bold role, the comment is the only italic one, and the
% two remaining roles differ in lightness as well as in hue.
\definecolor{kwblue}{RGB}{20,60,150}
\definecolor{cmtgreen}{RGB}{90,120,90}
\definecolor{strred}{RGB}{150,50,40}
\definecolor{numgray}{gray}{0.45}
\definecolor{framegray}{gray}{0.70}
\definecolor{bgshade}{gray}{0.97}
% A darker fill for figure cells. The listing background has to stay very
% light to keep code readable; a figure cell has to be clearly shaded or the
% distinction it draws is invisible in print.
\definecolor{cellshade}{gray}{0.90}

% The listing size is the document's choice, because it depends on the
% column: a two-column article sets `\listingsize` to 6pt and keeps its
% listings within 70 columns, a one-column document keeps the body size
% and includes source files at 80. Lines are never broken; a line that does not fit
% overflows the frame, which is visible, rather than wrapping, which is
% not.
\providecommand{\listingsize}{\scriptsize}

% `'` is deliberately not a string delimiter: the language uses it for loop
% labels, as Rust does, so treating it as a quote would swallow the rest of
% the line.
\lstdefinestyle{txhdl}{
  basicstyle=\ttfamily\listingsize,
  keywordstyle=\bfseries\color{kwblue},
  commentstyle=\itshape\color{cmtgreen},
  stringstyle=\color{strred},
  numberstyle=\tiny\color{numgray},
  morekeywords={mod,use,pub,const,type,struct,enum,trait,impl,fn,async,let,mut,
    match,if,else,loop,while,for,in,break,continue,return,as,await,
    module,bus,channel,transaction,protocol,pipeline,flow,seq,config,
    port,stage,pipe,instance,connect,bind,tag,domain,blackbox,foreign,
    property,role,signal,handler,true,false,
    sequence,interface,modport,implements,package,import,var,wire,func,proc,
    case,default,next,fork,branch,join,state,fsm},
  morecomment=[l]{//},
  morestring=[b]",
  showstringspaces=false,
  columns=fullflexible,
  keepspaces=true,
  breaklines=false,
  % Framed, so a listing is bounded rather than floating in the column.
  frame=single,
  rulecolor=\color{framegray},
  framesep=3pt,
  backgroundcolor=\color{bgshade},
  xleftmargin=3pt,
  xrightmargin=3pt,
  aboveskip=6pt,
  belowskip=6pt,
  % Every listing is numbered and captioned, so it can be referred to by
  % number. `caption` without `float` numbers the listing and leaves it
  % where it was written, which is what a listing in running prose needs.
  captionpos=b,
  abovecaptionskip=2pt,
  belowcaptionskip=6pt,
}

% Figures drawn as text. Framed the same way, and with the highlighting
% switched off, because a box-and-arrow diagram is not code and half its
% words turning blue reads as an accident. Setting a style adds keys rather
% than clearing the ones already set, so the three roles are neutralised by
% hand rather than by leaving them out.
\lstdefinestyle{ascii}{
  basicstyle=\ttfamily\listingsize,
  keywordstyle=\ttfamily,
  commentstyle=\ttfamily,
  stringstyle=\ttfamily,
  columns=fullflexible,
  keepspaces=true,
  frame=single,
  rulecolor=\color{framegray},
  framesep=3pt,
  backgroundcolor=\color{bgshade},
  xleftmargin=3pt,
  xrightmargin=3pt,
  aboveskip=6pt,
  belowskip=6pt,
}
\lstset{style=txhdl}

% House TikZ styles. `shorten` lives on the arrow styles rather than on each
% arrow, so every arrow in the document gets the gap, including ones added
% later. TikZ clips a path at a node's border, so without it an arrowhead
% butts against the box with no gap at all. Keep `node distance` well above
% twice the shorten, or the arrow shrinks to nothing.
\tikzset{
  box/.style={draw,rounded corners=2pt,align=center,inner sep=4pt,
              font=\footnotesize},
  boxf/.style={box,fill=bgshade},
  lbl/.style={font=\scriptsize\itshape,align=center},
  fl/.style={-{Stealth[length=4pt]},shorten >=2.5pt,shorten <=2.5pt},
  fld/.style={fl,dashed},
}


\usepackage{graphicx}

\InputIfFileExists{buildstamp.tex}{}{}
\providecommand{\buildstamp}{Draft build (outside the Bazel flow).}

% A range of a source file, by its begin/end markers.
\newcommand{\gpupart}[4]{%
\lstinputlisting[rangebeginprefix=//\ begin\{,rangebeginsuffix=\},%
rangeendprefix=//\ end\{,rangeendsuffix=\},includerangemarker=false,%
linerange=#2-#2,caption={#3},label={#4}]{gpu/razboj/src/#1}}

\title{Razboj: A Minimal GPU in TxHDL}

\author{Filip Filmar and Dragiša Janković%
\thanks{Both authors are independent. Filip Filmar is the
corresponding author, at f@hdlfactory.com; Dragiša Janković is
reached at linkedin.com/in/dragisa-jankovic.}%
\thanks{This is the tutorial on the first design that uses the AXI
link~\cite{axi} for something other than a test, for a reader who has
met TxHDL through the step-by-step tutorial~\cite{tutorial}. Every
listing is included from the project repository \code{HDL/txhdl} at
build time, the picture is the one Razboj drew when the build ran
it, and the waveform is that run's. The cover~\cite{cover} lists
every document.}%
\thanks{The authors used a large language model, Claude, as an
assistant in exploring the concepts and in writing and constructing
the documents and the programs; every commit of the repository says
so and carries the prompts that led to it, verbatim.}%
\thanks{\buildstamp}}

\markboth{Razboj: A Minimal GPU in TxHDL}{Filmar and Janković: Razboj: A Minimal GPU in TxHDL}

\begin{document}
\maketitle

\begin{abstract}
A rasteriser reads a display list and writes pixels into memory.
Razboj, the minimal GPU of this document, is one: a unit of about a
hundred lines that walks a box one pixel a
cycle, keeps three edge functions as it goes, and writes the pixels
that are inside the primitive as single-beat AXI bursts; the
framebuffer is a second unit, a memory behind the other end of the
link. The display list is in that memory too, and the rasteriser
reads it over the same link, so a program that writes the list is all
it takes to give Razboj work. Both are written in the lowered subset,
so both are netlists,
and the build simulates each netlist against the run that drew the
picture in this document, under \code{nvc} and under Verilator. The
rasteriser is also the first client of the AXI link that is hardware
rather than a testbench, and it is what that half of the link is for.
\end{abstract}

\section{What It Draws}
\label{sec:g-what}

Figure~\ref{fig:g-scene} is the demonstration, drawn by the
rasteriser, written into the framebuffer over AXI, read back out of
it and written to a file by the same program. It is sixty-four by
sixty-four pixels and eight entries of a display list: a clear, three
rectangles and four triangles.

\begin{figure}[htbp]
\centering
\includegraphics[width=0.62\columnwidth]{scene.png}
\caption{The demonstration, as Razboj drew it. Eight entries, four
thousand and ninety-six pixels, thirteen thousand cycles.}
\label{fig:g-scene}
\end{figure}

\gpupart{scene.rs}{scene}{The display list of the picture above, from
\code{gpu/razboj/src/scene.rs}, which drew Figure~\ref{fig:g-scene}.}{lst:g-scene}

\section{The Display List}
\label{sec:g-list}

A display list is a little instruction set, and its three entries
carry different things. A clear carries a colour and nothing else. A
rectangle carries a box. A triangle carries three vertices. That is
what an \code{enum} says in Rust, and it is what a program writes.

\gpupart{op.rs}{op}{The instruction set: \code{Op} as a program writes
it, \code{Kind} as the hardware reads it, and \code{Insn} as it goes
on the wire.}{lst:g-op}

\code{Insn} is one shape, as wide as the widest entry needs, because
every field a step of the rasteriser reads has to be at a fixed place:
the lowered subset slices a channel's word at constant offsets, and a
variable-width encoding would want a length field and a shifter in
front of everything else. A clear therefore leaves the box and the
vertices at zero and a rectangle leaves the vertices at zero, and the
kind says which of them mean anything. That is a fixed-width
instruction word, and it costs what one costs: the widest entry sets
the width, and the narrow ones pay for it.

Between the two is an assembler, and it does what a host can do once
per primitive rather than once per pixel.

\gpupart{op.rs}{encode}{\code{Op::encode}: the assembler. It clips the
box to the screen, winds the triangle so that its inside is where
every edge function is non-negative, and drops an entry that would
draw nothing.}{lst:g-encode}

The kind is read by the hardware in one place and decides two things.
A triangle is the entry whose edge functions are tested; a clear is
the entry whose box the rasteriser supplies, from the screen size in
its own type, since a clear is the whole screen by definition and
there is nothing for the instruction to say. A rectangle is the plain
case, its box read and its edges ignored. Those three lines are the
whole of the decoder:

\begin{lstlisting}[frame=none,numbers=none]
let covered = (self.kind.get() != Kind::Tri) | (n0 & n1 & n2);
let clearing = self.skind.get() == Kind::Clear;
let wx = mux(clearing, zero16, self.sx0.get().resize::<16>());
\end{lstlisting}

A wider instruction set would grow that decoder and not the walk,
which is the property worth having: a line, a circle or a sprite would
each add a kind and a way to say whether a pixel is inside, and the
one pixel a cycle, the three edge accumulators and the AXI writes
would stay as they are.

\section{The List in Memory}
\label{sec:g-mem}

\code{Insn} is how the list is held in a program and in a model.
In the memory it is words, and the format of those words is the one
thing software and hardware both have to agree on, so it is stated
once, in \code{gpu/razboj/src/dl.rs}.

\gpupart{dl.rs}{format}{The display list as a program writes it:
eight words an instruction, six of which say anything, and no field
across a word.}{lst:g-format}

An instruction takes eight words, so that the address of the
$i$th is the list's base plus $i$ shifted by five, a shift rather
than a multiply. No field crosses a word, so a program writes a field
with a shift and a mask and the rasteriser reads it with a slice at a
constant offset. Beside the list is one more word, the count. The
rasteriser reads the count until it is not zero, and then the list;
a program therefore writes the list first and the count last, and
writing the count is how it says the list is ready. Nothing else
passes between the two.

So the rasteriser has a front end. It polls the count, fetches the
six words of an instruction one read at a time, latching each field
as its word lands, draws the instruction, and fetches the next; when
the count is used up it is done and says so on \code{idle}. The
third vertex is not latched: its word is the last, so the setup reads
it straight off the bus in the cycle the walk starts. One read is out
at a time, because the front end has nothing to do until the answer
lands. The fetch costs the demonstration six reads an entry,
which is a few hundred cycles out of thirteen thousand.

\section{The Rasteriser}
\label{sec:g-raster}

An edge function is linear. For the edge from $a$ to $b$,
\[ e(x, y) = (b_x - a_x)(y - a_y) - (b_y - a_y)(x - a_x), \]
which is zero on the edge, positive on one side and negative on the
other. Stepping one pixel to the right adds $-(b_y - a_y)$ and
stepping one row down adds $(b_x - a_x)$, whatever the pixel. So the
rasteriser keeps, for each of the three edges, its value at the
current pixel, its value at the start of the current row, and those
two steps.

\gpupart{raster.rs}{state}{The rasteriser's state.}{lst:g-state}

The only multiplications are the six of the setup, done in the cycle
an entry is taken; the per-pixel work is three adds and three sign
tests. The arithmetic is thirty-two bits of two's complement, which
is wide enough for every product of two differences of vertices in
the range an entry may use. It was sixteen bits first, and a
pseudorandom scene with a vertex off the screen found the overflow
within a dozen tries, which is the argument for the soak test rather
than for the directed ones.

\gpupart{raster.rs}{run}{The rasteriser, two processes.}{lst:g-run}

The rasteriser is two processes. The first takes the link's answers
every cycle, says whether the pixel under the walk is in the
primitive, and counts what is in flight. The second is the front end
written as the sequence it is: poll the count until the list is
ready, then for each instruction fetch its six words, one read at a
time and each latched as it lands, and walk its box, every row and
every column a counted loop, a pixel a turn; when the list is drawn
it waits for ever. The lowering numbers the waits into a state
register of its own and gives each counted loop a counter, so the
netlist keeps what the hand-numbered machine before it kept in a
phase and a step without the source naming either.

The walk is one pixel per turn: the column advances and each edge
takes its column step, and at the end of a row the column goes back
to the first and each edge is reloaded from its row value plus its
row step, with one cycle between rows for the box to read back. A
pixel is written only when the link has room for the burst and its
beat, and the turn holds until then, so backpressure from the bus
stops the rasteriser rather than losing pixels.

\section{Writing Over AXI}
\label{sec:g-axi}

The rasteriser is an AXI host, and it is a host written as hardware.
The \code{Host} end of a link~\cite{axi} is a simulation-side
convenience, with an inbox, a table of answers and futures to await;
none of that lowers, and none of it needs to. What is underneath is
channels, and a client that is a unit holds those channel ends
itself: the issue channel and the write beats going out, the grant,
the write responses and the release coming back. \code{axi\_units}
hands them out under the names \code{HostClient} and
\code{PerClient}, and that is the whole of the difference between a
client that is a testbench and a client that is hardware.

A pixel is one burst of one beat: the issue and its beat go out in
the same cycle, so the write data channel stays in the order the
address channel is in, which is what AXI4 asks of a write. The
identifier the tracker granted is not needed, since a write's answer
matters only as a count and the front end has one read out at a time,
so it is taken and dropped. An identifier is handed back on
\code{release} as the answer arrives: a write's with its response, a
read's with its beat. That is what keeps four transactions in flight
and no more. A write's response is taken before a read's beat in the
cycle both arrive, since only one identifier goes back a cycle, and
the read waits a cycle for it.

\begin{figure*}[t]
\centering
\input{razboj_timing.tex}
\caption{The triangle fetched and begun. \code{word} counts the
instruction's words as the front end reads them, which land one at
a time on \code{rdata}; as the last word lands \code{kind} turns
from \code{Rect} to \code{Tri}. From
there \code{x} and \code{y} walk the triangle's box and \code{hit} is
whether the pixel is inside it; where \code{hit} is low no burst is
issued, so the row is walked and nothing is written. \code{inflight}
is the writes whose response has not come back.}
\label{fig:g-wave}
\end{figure*}

\section{The Framebuffer}
\label{sec:g-fb}

The framebuffer is the peripheral end of the same link: a memory of
one word per pixel, answering single-beat bursts.

\gpupart{fb.rs}{run}{The framebuffer.}{lst:g-fb}

It holds the display list as well as the pixels, above them. A read,
which is only ever the rasteriser fetching that list, is answered
from the memory in the cycle the request is taken.
A write is held until its beat arrives, because AXI4 puts no
identifier on the write data channel, so the beat belongs to the
oldest address phase and there is nothing to match it by but order.
One write is held at a time, which is why the demonstration takes
about three cycles a pixel rather than one; a framebuffer that took a
request and its beat together in the cycle both were offered would be
the first thing to improve.

\section{The Netlists}
\label{sec:g-netlist}

Both units lower. The traced run is sixteen by sixteen pixels, small
enough that a testbench generated from its trace replays in a moment,
and the build checks both netlists against it at their ports, as VHDL
under \code{nvc} and as Verilog under Verilator. Listing~\ref{lst:g-out}
is what that run printed.

The framebuffer's netlist is told what its memory starts with, the
list and its count, which \code{Mem::with} gave the run and which the
lowering cannot see; without it the simulated module reads zeroes
where the run read the list. The rasteriser is also the first unit
checked this way that does something in its first cycle, since it
goes looking for the count at once, and the testbench generator
starts a channel's \code{ready} high for that reason: an empty
channel has room, and the first cycle has no sample in the trace to
say so.

\lstinputlisting[style=ascii,caption={What \code{//gpu/razboj:trace} printed
when the build ran it: the same three kinds of entry on a screen
small enough to print.},label={lst:g-out}]{out_razboj.txt}

Three names had to be changed to lower. A register called \code{next}
gives VHDL a signal called \code{next}, which it reserves; a register
called \code{tri} and a wire called \code{inside} give Verilog two of
its own reserved words. The lowering emits a name as it is written,
and it did not then reject the reserved words of its targets, so the
failure was an analysis error in a generated file rather than an error
at the point the name was chosen. It now refuses such a name where the
name is declared. The registers here are \code{edged} and the wire is
\code{hit}.

A fourth mistake is worth stating because the lowering did not catch
it either: a \code{let} that computes, and takes the name of a
register, produces two declarations of one signal. That is now refused
too, at the \code{let}, once the unit's \code{lowered} is compiled.
The locals in the rasteriser are
\code{px}, \code{py} and \code{held} for that reason, where \code{x},
\code{y} and \code{waiting} would read better.

\section{What It Is Not}
\label{sec:g-not}

It is minimal, and the ways in which it is minimal are worth naming.
There is one pixel per cycle and no more; no depth buffer, so
primitives are drawn in the order the display list gives them; no
interpolation, so a triangle is one colour; no texture; no clipping
of geometry, only of the box; and one read out at a time, so the
fetch is as slow as the link's round trip. Every one of those is a
thing to add on top of what is here rather than a thing to work
around.

A CPU writing the display list is the step the bus was for, and it is
taken in the network document~\cite{noc}: the core~\cite{vreteno}
works out an icosahedron and writes it into a memory Razboj shares,
two hops away on a lattice.

\begin{thebibliography}{9}

\bibitem{axi}
F.~Filmar and D.~Janković, ``AXI in TxHDL,'' project repository
\code{HDL/txhdl}, \code{//docs:axi}, 2026.

\bibitem{tutorial}
F.~Filmar and D.~Janković, ``TxHDL, step by step,'' project
repository \code{HDL/txhdl}, \code{//docs:tutorial}, 2026.

\bibitem{vreteno}
F.~Filmar and D.~Janković, ``Vreteno: an RV32IMC core in TxHDL,''
project repository \code{HDL/txhdl}, \code{//docs:vreteno}, 2026.

\bibitem{noc}
F.~Filmar and D.~Janković, ``A network on chip in TxHDL,'' project
repository \code{HDL/txhdl}, \code{//docs:noc}, 2026.

\bibitem{cover}
F.~Filmar and D.~Janković, ``TxHDL documents,'' project repository
\code{HDL/txhdl}, \code{//docs:cover}, 2026.

\end{thebibliography}

\end{document}
